\section{The discrete logarithm problem}

\subsection{Formal definition}

\begin{definition}[Discrete logarithm problem]
Let $\G=\langle g\rangle$ be a cyclic group of order $r$. Given $g$ and
$h\in\G$, the discrete logarithm problem is to find
\[
  x\in\Z/r\Z\quad\text{such that}\quad h=g^x.
\]
In additive notation, given $G$ and $Q\in\langle G\rangle$, find $x$ such that
\[
  Q=[x]G.
\]
\end{definition}

The group, generator, and generator order are part of the problem statement.
The answer is unique only modulo $\ord(g)$. If the target is outside the
generated subgroup, no answer exists.

\subsection{What ``efficient'' and ``hard'' mean}

Write
\[
  n=\lceil\log_2 r\rceil,
\]
the bit length of the exponent space. Complexity must be measured in $n$, not
in the numerical value $r$. After reducing $x$ modulo $r$, binary
square-and-multiply uses $n-1$ squarings and at most $n-1$ additional
multiplications. Thus
\[
  x\longmapsto g^x
  \qquad\text{costs fewer than $2n$ group operations.}
\]
Double-and-add gives the same $O(n)$ count for $x\mapsto[x]G$. A group
operation is not a single hardware instruction: modular multiplication and
reduction have their own bit cost. If field elements have $k$ bits and
$M(k)$ is the bit cost of a $k$-bit multiplication, finite-field
square-and-multiply costs roughly $O(nM(k))$ up to reduction factors. This is
still polynomial in the encoded input length.

\begin{center}
\begin{tikzpicture}[
  node distance=53mm,
  state/.style={draw=black!25,rounded corners,fill=black!2,inner sep=6pt},
  >={Stealth[length=2mm]}
]
\node[state] (secret) {secret exponent $x$};
\node[state,right=of secret] (public) {public element $h=g^x$};
\draw[->,thick,cryptographgreen] (secret) --
  node[above,text=cryptographgreen] {forward: $O(n)$ group operations} (public);
\draw[->,thick,cryptographorange] (public) to[bend left=38]
  node[below,text=cryptographorange,align=center]
  {generic classical inverse: $\Theta(2^{n/2})$\\for constant success}
  (secret);
\end{tikzpicture}
\end{center}

The reverse arrow combines an upper bound with a restricted lower bound.
Baby-step--giant-step and Pollard rho use $O(\sqrt r)=O(2^{n/2})$ group
operations. More precisely, Shoup bounded the success probability of a
classical generic attack using $T$ oracle operations by $O(T^2/r)$ when $r$ is
prime \cite{shoup1997lower}. Constant success therefore requires
$\Omega(\sqrt r)$ operations, and the generic classical complexity is
$\Theta(2^{n/2})$.

\begin{center}
\begin{tabularx}{\textwidth}{lllX}
\toprule
Problem & Known work & Class & Status \\
\midrule
$x\mapsto g^x$ or $[x]G$ & $O(n)$ group operations & polynomial & proven algorithm \\
generic DLP/ECDLP & $\Theta(2^{n/2})$ & exponential & generic-model theorem \\
chosen concrete curve & best known: generic & exponential & assumption; no unconditional bound \\
quantum DLP/ECDLP & $\operatorname{poly}(n)$ & polynomial & Shor's algorithm \\
\bottomrule
\end{tabularx}
\end{center}

The word \emph{generic} is essential. The lower bound deliberately ignores the
internal representation of a group element, so it cannot rule out a shortcut
that exploits a particular encoding. No unconditional theorem says that
Curve25519 ECDLP requires exponential classical time. Its hardness is an
assumption supported by the generic lower bound, decades of cryptanalysis, and
the absence of a curve-specific shortcut.

Complexity-theoretic cryptography states that assumption over a family of
groups, and for random instances rather than one hand-picked puzzle:

\begin{definition}[Classical DLP hardness assumption]
Let $(\G_\lambda=\langle g_\lambda\rangle)$ be an efficiently represented
family of prime-order groups, with orders $r_\lambda$. DLP is hard for the
family if every classical probabilistic polynomial-time algorithm $\mathcal A$
satisfies
\[
  \Pr_{x\leftarrow\Z/r_\lambda\Z}
  \left[
    \mathcal A(\mathsf{par}_\lambda,g_\lambda,g_\lambda^x)
      \equiv x\pmod{r_\lambda}
  \right]
  \le \operatorname{negl}(\lambda).
\]
The probability includes the randomness of $x$ and of $\mathcal A$;
$\operatorname{negl}$ decreases faster than every inverse polynomial.
\end{definition}

Within one fixed cyclic group, DLP is random self-reducible. Given an arbitrary
target $h=g^x$, choose random $t$ and send $hg^t=g^{x+t}$ to a solver for
random targets; subtracting $t$ from its answer recovers $x$. This connects
random-target and worst-target instances inside that group, but it does not
prove that the selected group family is hard.

This is an average-case one-wayness assumption, not a proof that DLP is
NP-hard. Such NP-hardness is not known, and worst-case NP-hardness alone would
not provide the random-instance hardness needed in cryptography. Curve25519 is
one fixed parameter set, so its ``security bits'' are a concrete estimate of
the best known attack cost rather than an asymptotic theorem.

\subsection{Generic attacks}

\begin{center}
\begin{tabular}{lll}
\toprule
Method & Time & Storage \\
\midrule
exhaustive search & $O(r)=O(2^n)$ & $O(1)$ \\
baby-step--giant-step & $O(\sqrt r)=O(2^{n/2})$ & $O(\sqrt r)$ \\
Pollard rho & expected $O(\sqrt r)=O(2^{n/2})$ & $O(1)$ \\
Pohlig--Hellman & controlled by the largest prime factor of $r$ & varies \\
\bottomrule
\end{tabular}
\end{center}

Baby-step--giant-step writes $x=im+j$ with $m=\lceil\sqrt r\rceil$ and
looks for a collision
\[
  h(g^{-m})^i=g^j.
\]
Pollard rho replaces the stored table by a pseudorandom walk and a cycle-finding
method \cite{pollard1978rho}. The classical generic-group lower bound explains
why the square root repeatedly appears \cite{shoup1997lower}.

Pohlig--Hellman reduces a DLP of order
$r=\prod_i q_i^{e_i}$ to DLPs in the prime-power factors \cite{pohlig1978}. A secure group
therefore needs a large prime factor; a large ambient field by itself is not
enough.

Finite-field multiplicative groups also admit index-calculus-style algorithms
that exploit their representation and beat the generic square-root bound. For
prime fields, the number field sieve has heuristic subexponential running time
of the form \cite{gordon1993nfs}
\[
  L_p[1/3,c]
  =\exp\!\left((c+o(1))(\log p)^{1/3}(\log\log p)^{2/3}\right).
\]
This is subexponential in $p$ but still superpolynomial in the encoded input
length $\log p$. No analogous general subexponential method is known for a
properly chosen elliptic-curve group. This is why a 255-bit elliptic-curve
field can target a security level that would require a much larger
finite-field Diffie--Hellman modulus.

All of these classical comparisons have a quantum boundary. Shor's algorithm
solves DLP, including ECDLP, in time polynomial in the encoded input size on a
sufficiently capable fault-tolerant quantum computer \cite{shor1997}. The
hardness assumption is therefore explicitly classical, not post-quantum.

\subsection{DLP, CDH, and DDH are different problems}

Let $g$ generate a group of prime order $r$ and choose unknown
$a,b\in\Z/r\Z$.

\begin{definition}[Computational Diffie--Hellman]
The $\CDH$ problem maps
\[
  (g,g^a,g^b)\longmapsto g^{ab}.
\]
In additive notation it maps $(G,[a]G,[b]G)$ to $[ab]G$.
\end{definition}

\begin{definition}[Decisional Diffie--Hellman]
The $\DDH$ problem asks whether a supplied $T$ equals $g^{ab}$, given
$(g,g^a,g^b,T)$.
\end{definition}

A DLP solver gives a CDH solver: recover $a$, then compute $(g^b)^a$.
A CDH solver gives a DDH solver: compute $g^{ab}$ and compare it with $T$.
Thus there are solver implications
\[
  \DLP\Longrightarrow\CDH\Longrightarrow\DDH,
\]
so the implications between hardness assumptions run in the reverse direction:
\[
  \text{$\DDH$ hard}\Longrightarrow\text{$\CDH$ hard}
  \Longrightarrow\text{$\DLP$ hard}.
\]
These problems are not known to be equivalent in arbitrary groups. In
some pairing-friendly groups, for example, DDH is easy while CDH may remain
hard. Saying ``DLP is hard'' is therefore not itself a proof of every
Diffie--Hellman assumption.

\subsection{From one-wayness to cryptography}

For a prime-order group with a full point encoding, $x\mapsto g^x$ is a
bijection and a candidate \emph{one-way permutation}: everyone can evaluate
it, while inversion on a random output is assumed hard. This easy-forward,
hard-backward pattern is one of cryptography's recurring building blocks. Hash
functions seek preimage resistance, and trapdoor constructions add secret
information that enables inversion.

Diffie--Hellman uses the asymmetry differently: neither Alice nor Bob inverts
the public map. Each applies a scalar they already know to the other party's
public element. A passive outsider must compute the shared element without
knowing either scalar, which is the CDH problem. Thus one-wayness is necessary
intuition but not a complete security argument: key agreement needs a
CDH-style assumption and a full protocol still needs key derivation,
authentication, and message protection.

\begin{takeaway}
Curve25519 does not introduce a new kind of logarithm. It supplies a carefully
chosen cyclic subgroup in which the ordinary DLP is written additively and is
called ECDLP.
\end{takeaway}
